% Section 1: Introduction

\section{Introduction}
\label{sec:introduction}

Saturn’s F ring is one of the most dynamic and scientifically intriguing features in the solar system. Discovered in 1979 by the Pioneer 11 spacecraft \zcite{(Gehrels et al. 1980)}, the F ring is a narrow, complex, and ever-changing structure located just outside Saturn’s main ring system. Its unique morphology, characterized by multiple strands, clumps, and transient features, along with the influence of neighboring moons Prometheus and Pandora, has made it a focal point for studies of planetary ring dynamics, satellite-ring interactions, and the processes governing the evolution of ring systems.

Before 2004, quality images of the F ring were limited to a handful produced by the Voyager 1 and 2 spacecraft during their short fly-bys in 1980 and 1981. This all changed with the arrival of the Cassini Orbiter at Saturn on June 30, 2004. During its subsequent 13 years at Saturn, Cassini used its Imaging Science Subsystem (ISS) \zcite{(Knowles 2018; Porco et al. 2004)} to observe the F ring more than 100,000 times, providing a treasure trove of data that is unlikely to be surpassed for generations.

This \textit{User Guide} provides a detailed description of the organization and contents of the Cassini ISS F Ring Mosaics bundle archived at NASA’s Planetary Data System Ring-Moon Systems Node under NASA CDAP grant 80NSSC21K0527: “The recent history of Saturn's dusty rings”. The purpose of this archive is to provide a carefully curated and documented set of high-quality images of the F ring taken by the ISS Narrow Angle Camera (NAC) and Wide Angle Camera (WAC) \zcite{(Porco et al. 2004)} using clear filters. Special emphasis is placed on “FMOVIEs”, sequences of dozens to hundreds of images taken in rapid succession at a single inertial longitude while the F ring orbited through the field of view. Such movies provide a detailed picture of the F ring material at various corotating longitudes as it passes through some true anomaly in the F ring’s eccentric orbit.

All included images were calibrated in units of I/F using CISSCAL 4.0 \zcite{(Knowles 2018)} and then reprojected onto a grid with radius relative to the F ring core on the vertical axis and corotating longitude on the horizontal axis (see Section \ref{sec:reprojection}). Multiple images were stitched together to form more complete mosaics, and the mosaics were then analyzed to remove any background gradient.

The archive includes reprojected images, mosaics, and mosaics that have been adjusted by subtracting a background model, organized by Cassini “observation name” (see Section \ref{sec:image-selection}). Each reprojected image and mosaic is supplemented by a table containing metadata such as resolution, phase angle, and emission angle for each valid corotating longitude. Reprojected images have an additional file describing the spacecraft pointing used to perform the reprojection. Mosaics have an additional file listing the reprojected images that were used to create the mosaic. A series of browse images in various sizes is provided for each reprojected image and mosaic as well. In total 20,303 reprojected images are provided along with 305 associated mosaics.

Finally, three global index files are provided covering the complete set of reprojected images, mosaics, and background-subtracted mosaics. These index files summarize metadata such as time, corotating longitude coverage, lighting and viewing geometry, and orbit details for each data product, making it easy to search for and analyze the data products \textit{en masse}.

First time users are encouraged to read or skim this entire \textit{User Guide} at least once (and possibly twice) before making use of the archived results.
